\centerline{\bf \Huge Матбой}

\begin{flushright}
        \scriptsize{Нормальный человек синусами не занимается\\ Саша}
\end{flushright}

\problem{1}
На доске написаны числа $1^2, 2^2, \ldots, 100^2$. Найдите наименьшее натуральное число, на которое делятся ровно два из выписанных чисел.
%комол 8 класс утюм 1 задача

\problem{2}
В треугольнике $A B C$, в котором $\angle A=60^{\circ}$, проведена медиана $A M$. На продолжении стороны $A C$ за точку $C$ отмечена точка $N$ так, что $A M=M N$. Найдите $A B / C N$.
%комол 7 клас утюм 4 задача


\problem{3}
Можно ли на гранях восьми одинаковых кубиков записать числа (не обязательно натуральные) так, чтобы из этих кубиков можно было составить вдвое больший куб с любой суммой от 1 до 10000000 на видимых гранях?

\problem{4}
Число $x^2+xy+y^2$ заканчивается на $\ldots0$ Докажите, что тогда оно заканчивается на $\ldots00.$

\problem{5}
Положительные числа $x$, $y$, $z$ таковы, что $x+y+z=1/x+1/y+1/z$. Докажите, что $x+y+z\geq\sqrt{\dfrac{xy+1}{2}}+\sqrt{\dfrac{yz+1}{2}}+\sqrt{\dfrac{zx+1}{2}}$.


\problem{6}
Существует ли многочлен $f(x)$ четвертой степени с целыми коэффициентами такой, что для любого целого $k$ многочлен $f(x)+k$ не представляется в виде произведения двух квадратных трехчленов с целыми коэффициентами?


\problem{7}
На летние сборы ЭлМШ пришли $n$ учеников, некоторые из которых знакомы друг с другом. Оказалось, что на сборах у каждого ученика не более $2k$ знакомых, но у каждых двух учеников, не знакомых друг с другом, не менее $k$ общих знакомых. Докажите, что $n \leq 6k$.


\problem{8}
В остроугольном треугольнике $A B C$ высоты $A D, B E$ и $C F$ пересекаются в точке $H.$ $P$ и $Q$ --- проекции точек $A$ и $H$ на $E F$ соответственно. Пусть $R$ - точка пересечения $D P$ и $Q H$. Докажите, что $HQ = HR$.